% ============================================================
%  macros.tex —— 数学符号与常用命令
% ============================================================

% ---------- 数集 ----------
\newcommand{\R}{\mathbb{R}}
\newcommand{\N}{\mathbb{N}}
\newcommand{\Z}{\mathbb{Z}}
\newcommand{\Q}{\mathbb{Q}}
\newcommand{\C}{\mathbb{C}}
\newcommand{\Rstar}{\mathbb{R}^{*}}

% ---------- 集合 ----------
\newcommand{\set}[1]{\left\{ #1 \right\}}                    % 列举式集合
\newcommand{\setof}[2]{\left\{ #1 \mid #2 \right\}}          % 描述式集合
\newcommand{\setcomp}[2][U]{\complement_{#1} #2}             % 补集
\newcommand{\card}[1]{\left| #1 \right|}                     % 集合元素个数

% ---------- 绝对值、范数、取整 ----------
\newcommand{\abs}[1]{\left\lvert #1 \right\rvert}
\newcommand{\norm}[1]{\left\lVert #1 \right\rVert}

% ---------- 区间 ----------
\newcommand{\intcc}[2]{\left[ #1,\,#2 \right]}               % 闭区间
\newcommand{\intoo}[2]{\left( #1,\,#2 \right)}               % 开区间
\newcommand{\intco}[2]{\left[ #1,\,#2 \right)}               % 左闭右开
\newcommand{\intoc}[2]{\left( #1,\,#2 \right]}               % 左开右闭

% ---------- 角、三角形、向量 ----------
\newcommand{\degree}{\ensuremath{^{\circ}}}
\newcommand{\vect}[1]{\overrightarrow{#1}}
\newcommand{\vecb}[1]{\boldsymbol{#1}}
\newcommand{\para}{\mathrel{/\!/}}                           % 平行

% ---------- 微积分与常用记号 ----------
\newcommand{\diff}{\mathrm{d}}
\newcommand{\dx}{\,\mathrm{d}x}
\newcommand{\ee}{\mathrm{e}}
\newcommand{\ii}{\mathrm{i}}
\newcommand{\limn}{\lim_{n\to\infty}}
\newcommand{\limx}[1]{\lim_{x\to #1}}

% ---------- 排版辅助 ----------
\newcommand{\blank}[1][2cm]{\underline{\hspace{#1}}}         % 填空横线
\newcommand{\points}[1]{\hfill\normalfont\small（#1 分）}     % 分值
\newcommand{\key}[1]{\textbf{#1}}                            % 答案强调

% ---------- 插图辅助 ----------
% 单图：\fig[占版心宽度比例]{文件名(不含路径/扩展名也可)}{图题}{标签}
% 例：  \fig[0.6]{demo-parabola}{二次函数的图象}{fig:parabola}
\newcommand{\fig}[4][0.75]{%
  \begin{figure}[H]
    \centering
    \includegraphics[width=#1\linewidth]{#2}%
    \caption{#3}%
    \label{#4}%
  \end{figure}%
}

% 并排两图：\figpair[单图宽度比例]{左文件}{左说明}{右文件}{右说明}{总图题}{标签}
\newcommand{\figpair}[7][0.48]{%
  \begin{figure}[H]
    \centering
    \begin{minipage}[t]{#1\linewidth}
      \centering
      \includegraphics[width=\linewidth]{#2}\\[2pt]{\small #3}
    \end{minipage}\hfill
    \begin{minipage}[t]{#1\linewidth}
      \centering
      \includegraphics[width=\linewidth]{#4}\\[2pt]{\small #5}
    \end{minipage}
    \caption{#6}\label{#7}
  \end{figure}%
}
